% GENERATED FILE. Edit canonical lessons, then run npm run generate:books.
% canonical-source-hash: fnv1a64:5892d1a624732701
% canonical-lessons: SA-C312-pravakta, SA-C312-aropah, SA-C312-nirakaroti, SA-C312-anrtam, SA-C312-sandigdhah

\chapter{Claims and Denials}
\label{ch:sa-claims-and-denials}
\hlchaptermodality{312}

\begin{chapteropening}
\textbf{By the end of this chapter:} \emph{I can report what someone claimed and how it was denied.}

Complete the last lesson of chapter 312: I can report what someone claimed and how it was denied.
\end{chapteropening}

\section[\texorpdfstring{\sk{प्रवक्ता} (pravaktā)}{प्रवक्ता (pravaktā)}]{\sk{प्रवक्ता} (pravaktā) — a spokesman}

\label{lesson:SA-C312-pravakta}



\begin{quote}
Before the new one: say the Sanskrit for doubts, then the Sanskrit for a misunderstanding.
\end{quote}

\subsection*{You'll want to know: \sk{प्रवक्ता}}
\textbf{\sk{प्रवक्ता}} — \emph{pravaktā} — \textquotedblleft{}a spokesman\textquotedblright{}.

\begin{quote}
\sk{प्रवक्ता} \sk{सभायां} \sk{वदति।}
\end{quote}

\begin{quote}
\emph{pravaktā sabhāyāṃ vadati.}
\end{quote}

\begin{quote}
The spokesman speaks at the meeting.
\end{quote}

A masculine noun.

\begin{cousinweb}
\textbf{\sk{प्रवक्ता}} (\emph{pravaktā}) is built from \emph{pra}, \textquotedblleft{}forth\textquotedblright{}, and \emph{vaktṛ}, \textquotedblleft{}a speaker\textquotedblright{}, from the root √vac, \textquotedblleft{}to speak\textquotedblright{}; MW gives \emph{pravaktṛ} as one who tells or relates, an announcer, and a spokesman is its modern use (MW).
\end{cousinweb}

\begin{grammarlens}[title={What you've built}]
One more word for telling what happened.
\end{grammarlens}

\subsection*{Your turn}
\begin{itemize}
\raggedright
  \item \emph{Say it:} \emph{pravaktā}
  \item \emph{Say it:} \emph{pravaktā}, once more
  \item \emph{Say it:} \emph{pravaktā sabhāyāṃ vadati.} — The spokesman speaks at the meeting.
  \item \emph{Recall:} say the Sanskrit for doubts, then the Sanskrit for a misunderstanding, then say \emph{pravaktā} again
\end{itemize}

\begin{tcolorbox}[breakable,colback=teal!4,colframe=teal!35!black,title={Before you move on}]
What does \sk{प्रवक्ता} mean? (\textquotedblleft{}A spokesman\textquotedblright{}.) Say it once more.
\end{tcolorbox}
\section[\texorpdfstring{\sk{आरोपः} (āropaḥ)}{आरोपः (āropaḥ)}]{\sk{आरोपः} (āropaḥ) — an allegation}

\label{lesson:SA-C312-aropah}



\begin{quote}
Before the new one: say the Sanskrit for a misunderstanding, then the Sanskrit for a spokesman.
\end{quote}

\subsection*{You'll want to know: \sk{आरोपः}}
\textbf{\sk{आरोपः}} — \emph{āropaḥ} — \textquotedblleft{}an allegation\textquotedblright{}.

\begin{quote}
\sk{पत्रिका} \sk{न्यायाधीशे} \sk{आरोपं} \sk{प्रकाशयति।}
\end{quote}

\begin{quote}
\emph{patrikā nyāyādhīśe āropaṃ prakāśayati.}
\end{quote}

\begin{quote}
The newspaper publishes an allegation against the judge.
\end{quote}

A masculine noun.

\begin{cousinweb}
\textbf{\sk{आरोपः}} (\emph{āropaḥ}) is built from \emph{ā} and the causative of the root √ruh, \textquotedblleft{}to rise\textquotedblright{}, as in \emph{āropayati}; MW gives \emph{āropa} as imposing a burden, charging with, attributing to, and an allegation is its modern use (MW).
\end{cousinweb}

\begin{grammarlens}[title={What you've built}]
One more word for telling what happened.
\end{grammarlens}

\subsection*{Your turn}
\begin{itemize}
\raggedright
  \item \emph{Say it:} \emph{āropaḥ}
  \item \emph{Say it:} \emph{āropaḥ}, once more
  \item \emph{Say it:} \emph{patrikā nyāyādhīśe āropaṃ prakāśayati.} — The newspaper publishes an allegation against the judge.
  \item \emph{Recall:} say the Sanskrit for a misunderstanding, then the Sanskrit for a spokesman, then say \emph{āropaḥ} again
\end{itemize}

\begin{tcolorbox}[breakable,colback=teal!4,colframe=teal!35!black,title={Before you move on}]
What does \sk{आरोपः} mean? (\textquotedblleft{}An allegation\textquotedblright{}.) Say it once more.
\end{tcolorbox}
\section[\texorpdfstring{\sk{निराकरोति} (nirākaroti)}{निराकरोति (nirākaroti)}]{\sk{निराकरोति} (nirākaroti) — denies; rejects}

\label{lesson:SA-C312-nirakaroti}



\begin{quote}
Before the new one: say the Sanskrit for a spokesman, then the Sanskrit for an allegation.
\end{quote}

\subsection*{You'll want to know: \sk{निराकरोति}}
\textbf{\sk{निराकरोति}} — \emph{nirākaroti} — \textquotedblleft{}denies; rejects\textquotedblright{}.

\begin{quote}
\sk{न्यायाधीशः} \sk{आरोपं} \sk{दृढं} \sk{निराकरोति।}
\end{quote}

\begin{quote}
\emph{nyāyādhīśaḥ āropaṃ dṛḍhaṃ nirākaroti.}
\end{quote}

\begin{quote}
The judge firmly denies the allegation.
\end{quote}

Past (he …): \sk{सः} \sk{आरोपं} \sk{निराकृतवान्}, \emph{saḥ āropaṃ nirākṛtavān}.

\begin{cousinweb}
\textbf{\sk{निराकरोति}} (\emph{nirākaroti}) is built from \emph{nir}, \emph{ā} and the root √kṛ, \textquotedblleft{}to do\textquotedblright{}; MW gives \emph{nir-ā-kṛ} as to drive away, repudiate, reject and contradict (MW).
\end{cousinweb}

\begin{grammarlens}[title={What you've built}]
One more word for telling what happened.
\end{grammarlens}

\subsection*{Your turn}
\begin{itemize}
\raggedright
  \item \emph{Say it:} \emph{nirākaroti}
  \item \emph{Say it:} \emph{nirākaroti}, once more
  \item \emph{Say it:} \emph{nyāyādhīśaḥ āropaṃ dṛḍhaṃ nirākaroti.} — The judge firmly denies the allegation.
  \item \emph{Recall:} say the Sanskrit for a spokesman, then the Sanskrit for an allegation, then say \emph{nirākaroti} again
\end{itemize}

\begin{tcolorbox}[breakable,colback=teal!4,colframe=teal!35!black,title={Before you move on}]
What does \sk{निराकरोति} mean? (\textquotedblleft{}Denies; rejects\textquotedblright{}.) Say it once more.
\end{tcolorbox}
\section[\texorpdfstring{\sk{अनृतम्} (anṛtam)}{अनृतम् (anṛtam)}]{\sk{अनृतम्} (anṛtam) — a lie; an untruth}

\label{lesson:SA-C312-anrtam}



\begin{quote}
Before the new one: say the Sanskrit for an allegation, then the Sanskrit for denies.
\end{quote}

\subsection*{You'll want to know: \sk{अनृतम्}}
\textbf{\sk{अनृतम्}} — \emph{anṛtam} — \textquotedblleft{}a lie; an untruth\textquotedblright{}.

\begin{quote}
\sk{सः} \sk{आरोपः} \sk{अनृतम्} \sk{अस्ति।}
\end{quote}

\begin{quote}
\emph{saḥ āropaḥ anṛtam asti.}
\end{quote}

\begin{quote}
That allegation is a lie.
\end{quote}

A neuter noun.

\begin{cousinweb}
\textbf{\sk{अनृतम्}} (\emph{anṛtam}) is built from \emph{an}, \textquotedblleft{}not\textquotedblright{}, and \emph{ṛta}, \textquotedblleft{}right, true\textquotedblright{}; MW gives \emph{anṛta} as not true and, as a noun, falsehood, lying (MW).
\end{cousinweb}

\begin{grammarlens}[title={What you've built}]
One more word for telling what happened.
\end{grammarlens}

\subsection*{Your turn}
\begin{itemize}
\raggedright
  \item \emph{Say it:} \emph{anṛtam}
  \item \emph{Say it:} \emph{anṛtam}, once more
  \item \emph{Say it:} \emph{saḥ āropaḥ anṛtam asti.} — That allegation is a lie.
  \item \emph{Recall:} say the Sanskrit for an allegation, then the Sanskrit for denies, then say \emph{anṛtam} again
\end{itemize}

\begin{tcolorbox}[breakable,colback=teal!4,colframe=teal!35!black,title={Before you move on}]
What does \sk{अनृतम्} mean? (\textquotedblleft{}A lie; an untruth\textquotedblright{}.) Say it once more.
\end{tcolorbox}
\section[\texorpdfstring{\sk{सन्दिग्धः} (sandigdhaḥ)}{सन्दिग्धः (sandigdhaḥ)}]{\sk{सन्दिग्धः} (sandigdhaḥ) — doubtful; questionable}

\label{lesson:SA-C312-sandigdhah}



\begin{quote}
Before the new one: say the Sanskrit for denies, then the Sanskrit for a lie.
\end{quote}

\subsection*{You'll want to know: \sk{सन्दिग्धः}}
\textbf{\sk{सन्दिग्धः}} — \emph{sandigdhaḥ} — \textquotedblleft{}doubtful; questionable\textquotedblright{}.

\begin{quote}
\sk{प्रवक्तुः} \sk{उत्तरं} \sk{सन्दिग्धम्} \sk{अस्ति।}
\end{quote}

\begin{quote}
\emph{pravaktuḥ uttaraṃ sandigdham asti.}
\end{quote}

\begin{quote}
The spokesman's answer is doubtful.
\end{quote}

\begin{cousinweb}
\textbf{\sk{सन्दिग्धः}} (\emph{sandigdhaḥ}) is built from \emph{sam} and the root √dih, like \emph{sandeha}; MW gives \emph{saṃdigdha} as smeared over, confused, and questionable (MW).
\end{cousinweb}

\begin{grammarlens}[title={What you've built}]
One more word for telling what happened.
\end{grammarlens}

\subsection*{Your turn}
\begin{itemize}
\raggedright
  \item \emph{Say it:} \emph{sandigdhaḥ}
  \item \emph{Say it:} \emph{sandigdhaḥ}, once more
  \item \emph{Say it:} \emph{pravaktuḥ uttaraṃ sandigdham asti.} — The spokesman's answer is doubtful.
  \item \emph{Recall:} say the Sanskrit for denies, then the Sanskrit for a lie, then say \emph{sandigdhaḥ} again
\end{itemize}

\begin{tcolorbox}[breakable,colback=teal!4,colframe=teal!35!black,title={Before you move on}]
What does \sk{सन्दिग्धः} mean? (\textquotedblleft{}Doubtful; questionable\textquotedblright{}.) Say it once more.
\end{tcolorbox}
